% Modernised reading — fonds Alexandre Grothendieck, shelfmark 152, pages 1-12.
%
% This is an INTERPRETATION, not a transcription. Everything the critical
% apparatus carried moves into footnotes. In any dispute of substance the
% transcription governs: whoever cites Grothendieck must cite batch-01.fr.tex.
%
% Pass: Opus 5 (claude-opus-5), 2026-09-13 — first pass, unchecked against the
% pages by a human. The folder was read whole; it holds one batch, pages 1-12.
% Pages 5, 8, 11 are listings with nothing of his.
%
% What the folder is. Pages 1-3 are trial sheets on « types de mots » and the
% chord diagrams that code them, with a recursive notion of admissibility.
% Pages 4-12 (his 1, 1', 2, 2', 3, 3') are one text: a plane tree is a tree
% with a rotation system up to reversal; cutting the surface along a
% circularised graph gives a contour; circularised graphs are equivalent to
% oriented contours with a fixed-point-free involution on their edges; a tree
% is exactly one polygon with a non-crossing involution, i.e. a disc cut by
% non-crossing chords, whose dual graph is the tree.
%
% Notation fixed for the whole document. Gamma a graph, S its vertices, A its
% edges, A_s the edges at s, A-bar the arcs (darts, oriented edges),
% sigma_Gamma the orientation reversal, rho_Gamma the rotation (next dart in
% the circular order at the origin). circ(Gamma) the set of circularisations.
% C the contour, A(C) its edges, sigma_C the gluing involution on A(C), rho_C
% the successor along the oriented contour. rd, rg: right-hand and left-hand
% edge of a dart (his abbreviations).
%
% Substantive departures from the pages, each footnoted where it occurs:
% — page 4: « si les sommets de Gamma sont d'ordre <= 1, il y a une seule »
%   circularisation; the right bound is <= 2 (cyclic orders on sets of at
%   most two elements are unique), as the preceding sentence (card A_s >= 3)
%   says.
% — pages 7 and 9: the dictionary is stated three ways on the leaves and they
%   disagree: page 7's box rho_C(rd a) = rd(rho_Gamma a), page 9's head
%   sigma_Gamma = sigma_C and rho_Gamma = rho_C, and page 9's table, whose
%   last two rows give rho_C <-> rho_Gamma sigma_Gamma and rho_Gamma <->
%   rho_C sigma_C. The consistent dictionary is sigma_C = sigma_Gamma,
%   rho_C = rho_Gamma sigma_Gamma (face permutation), hence rho_Gamma =
%   rho_C sigma_C; the reading states it and footnotes the three variants.
% — page 2: the correspondence « types de mots » <-> chord diagrams is given
%   in the form the leaf states; the reading of what a « type de mot » is
%   (cyclic words in which letters are cancelled in pairs) is ours and marked.
%
% Modern names given and footnoted as ours: rotation system (Heffter,
% Edmonds), ribbon graph, combinatorial map (darts with an involution and a
% rotation, faces as cycles of their product), chord diagram, non-crossing
% matching, dual tree of a dissection, Catalan numbers for rooted plane trees,
% free reduction.
%
% What the folder announces and never establishes: the equivalence of page 2
% (stated); the two last words of the boxed criterion of page 9 (illegible);
% the functoriality of the equivalences, asserted as « équivalences de
% catégories » and not checked on morphisms.
\documentclass[11pt,a4paper]{article}
% Le préambule commun à toutes les transcriptions du fonds Grothendieck.
%
% Il tient en une idée : **l'incertitude se compose, elle ne se résout pas.**
% Une transcription qui lisse un mot douteux a détruit la seule chose pour
% laquelle on la faisait — un lecteur doit pouvoir savoir, à chaque endroit,
% ce qui était sur le feuillet et ce qui a été ajouté. Les six macros qui
% suivent sont donc l'appareil critique, et non de la décoration.
%
% Les mêmes commandes sont reconnues par `scripts/render.mjs`, qui produit la
% vue de lecture du site. Une macro ajoutée ici doit l'être aussi là-bas,
% faute de quoi le rendu échouera bruyamment — ce qui est voulu : un rendu
% silencieusement faux est pire qu'un rendu absent.

\ProvidesPackage{grothendieck}[2026/08/08 Transcriptions du fonds Grothendieck]

\usepackage{amsmath,amssymb,amsthm}
\usepackage{mathrsfs}
\usepackage{stmaryrd}
% Les diagrammes commutatifs sont partout dans ce fonds : c'est la langue
% même de la géométrie algébrique de Grothendieck, et une transcription qui
% les rendrait en prose ne transcrirait rien.
\usepackage{tikz-cd}
% Français par défaut — c'est la langue des feuillets — mais l'anglais est
% chargé aussi : la traduction et le résumé sont des documents anglais, et la
% typographie française (espaces avant les deux-points) y serait une faute.
% Ces documents basculent par \selectlanguage{english} en tête de corps.
\usepackage[english,french]{babel}
\usepackage{fontspec}
\usepackage{geometry}
\geometry{a4paper,margin=25mm,marginparwidth=28mm}
\usepackage{marginnote}
\usepackage{xcolor}
% ulem, not soul. `soul`'s \st and \ul are built on a character-by-character
% scan that XeLaTeX's font machinery breaks: the document fails with an
% undefined control sequence rather than degrading. ulem does the same two jobs
% and is Unicode-engine safe. `normalem` stops it hijacking \emph, which the
% transcriptions use for Grothendieck's own emphasis.
\usepackage[normalem]{ulem}
\usepackage{hyperref}
\hypersetup{colorlinks=true,linkcolor=black,urlcolor=black,pdfborder={0 0 0}}

% --- Métadonnées du lot -----------------------------------------------------
% Reprises telles quelles par le script de rendu, qui les affiche en tête de
% la vue de lecture. Elles fixent ce que le fichier prétend couvrir : sans
% elles, une transcription flotte sans point d'attache dans un fonds de seize
% mille feuillets.

\newcommand{\folder}[1]{\gdef\grothFolder{#1}}
\newcommand{\batch}[1]{\gdef\grothBatch{#1}}
\newcommand{\leaves}[2]{\gdef\grothFirst{#1}\gdef\grothLast{#2}}
\newcommand{\foldertitle}[1]{\gdef\grothFoldertitle{#1}}
\gdef\grothFolder{?}\gdef\grothBatch{?}\gdef\grothFirst{?}\gdef\grothLast{?}\gdef\grothFoldertitle{}

% --- Le résumé --------------------------------------------------------------
%
% Il ouvre la lecture modernisée, sous le titre « Résumé », et il est composé
% en petit. Ce n'est pas un ornement : c'est le seul passage du document qui
% ne soit pas des mathématiques, et un lecteur doit pouvoir le reconnaître
% avant de l'avoir lu — puis le sauter s'il connaît déjà le sujet, ou s'y
% arrêter s'il ne le connaît pas. Le filet qui le suit marque la reprise du
% corps.

\newenvironment{resume}
  {\small\setlength{\parskip}{0.45em}}
  {\par\medskip\hrule\medskip}

% --- Mots-clés ---------------------------------------------------------------
%
% En anglais, en clôture du résumé de la lecture modernisée : le vocabulaire
% moderne sous lequel chercher ce que les feuillets construisent. Le manifeste
% du site (`npm run manifest`) les extrait pour étiqueter la cote entière —
% cette ligne est la source unique des tags, il n'y a pas de fichier à côté.
\newcommand{\keywords}[1]{\par\smallskip\noindent
  {\footnotesize\textsc{Keywords} --- \textit{#1}}\par}

% --- L'appareil critique ----------------------------------------------------

% Le numéro de feuillet, en marge. C'est l'ancre qui permet de régler un
% désaccord sur une formule en regardant *un* feuillet plutôt qu'une chemise
% de six cents. Le site s'en sert aussi pour tourner les pages du fac-similé
% quand on fait défiler la transcription.
\newcommand{\leaf}[1]{%
  \marginnote{\footnotesize\color{gray!70}\textsf{#1}}%
  \label{leaf:#1}\ignorespaces}

% Illisible : on ne devine pas. Un mot inventé qui se lit comme les autres est
% le pire résultat possible.
\newcommand{\ill}{\textcolor{red!60!black}{[\,\ldots\,]}}

% Lecture douteuse : le mot est donné, et signalé comme incertain.
\newcommand{\uncertain}[1]{\uline{#1}}

% Ajout de l'éditeur — une lettre manquante, un mot sous-entendu.
\newcommand{\add}[1]{\textcolor{blue!50!black}{[#1]}}

% Biffé par Grothendieck lui-même. Ce qu'il a rayé fait partie du document :
% une rature dit souvent où la pensée a changé de direction.
\newcommand{\struck}[1]{\sout{#1}}

% Note du transcripteur, sur l'état matériel du feuillet.
\newcommand{\note}[1]{\par\smallskip{\small\color{gray!60!black}\textit{[#1]}}\par\smallskip}

% Note marginale de Grothendieck, distincte du corps du texte.
\newcommand{\marginal}[1]{\par\smallskip\noindent\hspace{1em}%
  {\small\color{gray!60!black}#1}\par\smallskip}

% --- Titre ------------------------------------------------------------------

\AtBeginDocument{%
  \begin{center}
    {\large\bfseries Fonds Alexandre Grothendieck --- cote n\textsuperscript{o}~\grothFolder}\\[2pt]
    {\itshape\grothFoldertitle}\\[4pt]
    {\small feuillets \grothFirst--\grothLast\ (lot \grothBatch)}\\[2pt]
    {\footnotesize\color{gray!60!black}Transcription établie d'après le fac-similé numérisé
     par l'Université de Montpellier.\\
     Les lectures incertaines sont soulignées, les illisibles notées [\,\ldots\,],
     les ajouts entre crochets.}
  \end{center}
  \vspace{1em}}

\endinput


\folder{152}
\pages{1}{12}
\dating{[à partir de 1982]}
\watermark{Édition de démonstration\\interprétation personnelle de l'œuvre}
\foldertitle{Catégorie des types de mots [dont arbres plans] : notes manuscrites (s.d.) — lecture modernisée du dossier entier}

\begin{document}

\section*{Résumé}

\begin{resume}
Un arbre est un réseau de points reliés par des segments, sans aucune boucle.
Dessiné sur une feuille, il acquiert une information de plus : autour de chaque
point, l'ordre dans lequel les segments en partent, en tournant. Deux dessins
du même arbre peuvent différer par cet ordre, et c'est cette information qui
fait d'un arbre un \emph{arbre plan}. Ces feuillets en donnent trois
descriptions, et montrent qu'elles reviennent au même.

La première est celle qu'on vient de dire : un ordre circulaire autour de chaque
sommet, défini à renversement global près parce que la feuille n'a pas de
côté privilégié. La deuxième s'obtient en découpant la feuille le long de
l'arbre, comme on déplie un patron : on obtient un polygone, dont chaque côté
longe l'un des deux bords d'un segment, et une règle qui dit quels côtés
recoller. La troisième est le dessin dual : un disque découpé par des cordes qui
ne se croisent pas, chaque région devenant un sommet et chaque corde une arête.

Le passage de la première à la deuxième marche pour n'importe quel graphe
dessiné sur une surface, et non seulement pour les arbres ; on obtient alors
plusieurs polygones, un par face. C'est le dictionnaire des \emph{cartes
combinatoires} : on ne garde que les demi-arêtes, et deux permutations --- l'une
qui retourne une arête, l'autre qui tourne autour d'un sommet --- dont le
produit fait le tour des faces. L'arbre se reconnaît alors à ce qu'il n'a
qu'une face et que les recollements de son polygone ne se croisent pas.

Les trois premiers feuillets, des feuilles d'essais, attaquent le même objet
par un autre côté : des « types de mots » --- des suites de lettres et de leurs
inverses, à ce qu'il semble --- que des cordes dans un disque codent aussi. Le texte est celui d'une rédaction qui cherche ses notations : trois
formulations du dictionnaire y coexistent, et elles ne s'accordent pas toutes.

\keywords{plane tree, rotation system, ribbon graph, combinatorial map, chord diagram, non-crossing matching, polygon dissection, dual tree}
\end{resume}

\section*{Le fil du dossier, et les conventions}

\begin{enumerate}
\item[1.] \emph{Feuilles d'essais} (pages 1 à 3) : « types de mots », disques
découpés par cordes, admissibilité.
\item[2.] \emph{Arbres plans et circularisations} (page 4).
\item[3.] \emph{Découpage et contours} (page 6).
\item[4.] \emph{Le dictionnaire graphes circularisés / contours} (pages 7 et 9),
et le critère pour qu'un graphe soit un arbre.
\item[5.] \emph{Disques découpés par cordes} (pages 10 et 12).
\end{enumerate}

Les pages 4 à 12 portent sa pagination $1, 1', 2, 2', 3, 3'$ et forment un
texte suivi ; les pages 5, 8 et 11 sont des listings sans rien de sa main.

\emph{Conventions.} Un graphe $\Gamma$ a un ensemble de sommets $S$ et un
ensemble d'arêtes $A$ ; $A_{s}$ est l'ensemble des arêtes incidentes à $s$.
$\bar{A}$ est l'ensemble des \emph{arcs}, ou arêtes orientées ;
$\sigma_{\Gamma}$ est l'involution de $\bar{A}$ qui renverse l'orientation, et,
pour un graphe circularisé, $\rho_{\Gamma}$ la permutation de $\bar{A}$ qui envoie
un arc sur l'arc suivant, dans l'ordre circulaire, autour de son origine. Pour
un contour $C$, $A(C)$ est l'ensemble de ses arêtes, $\rho_{C}$ la permutation
« arête suivante » sur le contour orienté, et $\sigma_{C}$ l'involution de
recollement. $\mathrm{rd}(\bar{a})$ et $\mathrm{rg}(\bar{a})$ sont l'arête du
contour qui longe l'arc $\bar{a}$ à sa droite et à sa gauche.\footnote{« rd »
et « rg » sont ses abréviations ; la page 7 écrit la première en toutes lettres,
« arête droite de $\bar{a}$ ». $\bar{A}$ est sa notation pour l'ensemble des
arcs, employée sans définition à partir de la page 6.}

\emph{Ce que le dossier annonce et n'établit pas} : l'équivalence de la page 2,
énoncée ; la fin du critère encadré de la page 9, illisible ; et la
fonctorialité des équivalences, affirmées « de catégories » sans que les
morphismes soient examinés.

\pagerange{1}{3}
\section*{Feuilles d'essais : types de mots et cordes (pages 1 à 3)}

\emph{Page 1.} Un champ d'une trentaine de petites figures --- cercles marqués
de points et de rayons, arbres à trois ou quatre branches --- où l'on lit des
« systèmes » de deux lettres $uv$, $u^{-1}v$, $uv^{-1}$, $u^{-1}v^{-1}$, de
trois lettres $uvw$, les systèmes $u^{-1}u$ et $uu^{-1}$, des « sommets
principaux », des mots $uvu^{-1}$, et les mots « conjugués » et
« invariants ».

\emph{Page 2.} L'énoncé que le dossier se donne :

\textbf{Énoncé.} La catégorie des « types de mots » de longueur $n$ est
équivalente à la catégorie isotypique des systèmes $(D, \omega, K, S, S_{0})$
formés d'un disque $D$ orienté par $\omega$, d'un découpage de $D$ par un
ensemble $K$ de cordes, et d'un ensemble $S$ de $n$ sommets marqués sur le bord,
avec au moins un sommet sur chaque arc que les extrémités des cordes découpent
sur $\partial D$. Son \emph{nombre d'ordre réduit} est
\[
\sum_{\alpha \in \pi_{0}(\mathring{D} \smallsetminus K)}
\bigl(\mathrm{card}(\alpha \cap S) - 1\bigr),
\]
la somme portant sur les régions du découpage. Un type de mot est
\emph{trivial} quand ce nombre est nul.\footnote{L'énoncé ne définit pas
« type de mot », ni $S_{0}$, et la page ajoute « type de mot réduit :
$K = \emptyset$ (i.e.\ trivial) », qui ne s'accorde pas avec la définition du
trivial par l'ordre réduit nul. Un tableau en tête du feuillet aligne « disques
découpés par cordes $\leftrightarrow$ arbres, graphes circularisés
$\leftrightarrow$ contours $\simeq$ involutions » : c'est le programme des pages
4 à 12. L'adjectif « trivial » et le mot « contour » sont lus avec doute.}

\emph{Page 3.} Un graphe muni d'une involution $\sigma$ sur l'ensemble de ses
arêtes est dit \emph{admissible} par récurrence sur le nombre d'arêtes : si
$\sigma = \mathrm{id}$, ou s'il existe une orbite de $\sigma$ formée de deux arêtes
consécutives dont la contraction en un point donne un graphe admissible, muni de
l'involution induite. Tout graphe sans arête est admissible.\footnote{L'adjectif
qui qualifie le graphe est lu « segmenté » avec doute, et « admissible » aussi.
Si l'on pense les arêtes consécutives d'un contour comme les lettres d'un mot
cyclique, et les orbites de $\sigma$ comme les paires $x$, $x^{-1}$, la
récurrence est celle de la réduction libre --- on efface deux lettres voisines
inverses l'une de l'autre --- et un appariement est admissible exactement quand
ses cordes ne se croisent pas. Ce rapprochement est le nôtre ; il éclaire la
page 2 et le critère de la page 9, mais le feuillet ne le formule pas.}

\pagerange{4}{4}
\section*{Arbres plans et circularisations (page 4)}

Un \emph{arbre} est un graphe connexe et simplement connexe ; il est donc
strict --- sans boucle, et au plus une arête entre deux sommets distincts --- et
se décrit par $(S, A)$, $A \subset \mathfrak{P}_{2}(S)$, avec $S \neq \emptyset$,
et $A = \emptyset$ si et seulement si l'arbre est ponctuel.

\textbf{Définition.} Un \emph{arbre plan} est un arbre muni d'une classe de
plongements $|\Gamma| \hookrightarrow X$ dans une surface, deux plongements
dans $X$ et $X'$ étant identifiés s'il existe des voisinages $U$, $U'$ de
$|\Gamma|$ et un homéomorphisme $U \simeq U'$ induisant l'identité sur
$|\Gamma|$.

Une \emph{circularisation} de $\Gamma$ est la donnée, pour chaque sommet $s$,
d'un ordre circulaire sur $A_{s}$ ; $\mathrm{circ}(\Gamma)$ est l'ensemble des
circularisations, muni du passage à l'ordre opposé.

\textbf{Proposition.}
\begin{enumerate}
\item[(i)] Les classes de plongements de $|\Gamma|$ dans une surface
\emph{orientée} correspondent aux circularisations de $\Gamma$.
\item[(ii)] Les classes de plongements dans une surface quelconque correspondent
aux \emph{bicircularisations} : couples $(\varepsilon, c)$ formés d'un ensemble
$\varepsilon$ à deux éléments et d'une application $c : \varepsilon \to
\mathrm{circ}(\Gamma)$ qui échange les deux éléments de $\varepsilon$ avec deux
circularisations opposées. Autrement dit, aux circularisations prises à
renversement global près.
\end{enumerate}
En effet $|\Gamma|$ étant simplement connexe, il admet un voisinage simplement
connexe $U$, et $U \simeq \mathbb{R}^{2}$ ; les deux orientations de $U$
forment $\varepsilon$.\footnote{Le nom moderne d'une circularisation est
\emph{système de rotation} (L.~Heffter, 1891 ; J.~Edmonds, 1960), et un graphe
qui en est muni est un \emph{graphe à rubans} ; les noms ne sont pas sur la
page. Le mot qui introduit $(\varepsilon, c)$, lu « paires », suit le sens et non
le tracé.}

La circularisation ne dépend que des $A_{s}$ de cardinal au moins 3. Si tous
les sommets sont de degré au plus 2 --- l'arbre est un chemin ---, il n'y a
qu'une circularisation, et un seul arbre plan.\footnote{La page écrit « si les
sommets de $\Gamma$ sont d'ordre $\leq 1$ ». Un ordre circulaire sur un ensemble
à au plus deux éléments est unique, et la phrase précédente de la page, qui
parle des $A_{s}$ de cardinal $\geq 3$, demande la borne 2 ; on corrige.}

\pagerange{6}{6}
\section*{Découper la surface le long du graphe (page 6)}

En découpant la surface $X$ le long de $|\Gamma|$, $\Gamma$ étant un arbre, on
obtient une surface à bord connexe, munie d'une structure naturelle de polygone
topologique ; le polygone combinatoire correspondant ne dépend que de l'arbre
plan, à isomorphisme canonique près. Il a deux côtés par arête de l'arbre.

Pour un graphe circularisé quelconque, la même construction donne un
\emph{contour} --- une réunion disjointe finie de polygones, un par face ---
qui n'est plus connexe en général. Les sommets de $\Gamma$ se retrouvent sur les
arcs :
\[
S \simeq \bar{A}/\rho_{\Gamma}
\]
si $\Gamma$ n'a pas de sommet isolé, chaque orbite de la rotation étant
l'ensemble des arcs issus d'un sommet.\footnote{La seconde moitié du feuillet
est un bloc annulé par deux horizontales et trois diagonales ; on n'en garde que
la formule $S \simeq \bar{A}/\rho_{\Gamma}$, qui y est lisible et que la suite
utilise, et l'introduction de $\rho_{\Gamma}$ qui la précède.}

\pagerange{7}{9}
\section*{Graphes circularisés et contours : le dictionnaire (pages 7 et 9)}

\emph{Les arêtes du contour sont les arcs.} Chaque arc $\bar{a}$ est longé à sa
droite par une arête du contour, et l'application
\[
\mathrm{rd} : \bar{A} \xrightarrow{\;\sim\;} A(C),
\qquad \bar{a} \longmapsto \text{l'arête droite de } \bar{a},
\]
est bijective ; l'arête gauche est $\mathrm{rg}(\bar{a}) =
\mathrm{rd}(\sigma_{\Gamma}\bar{a})$. Par elle, chaque arête de $C$ reçoit une
orientation naturelle, et $C$ est orienté.

\emph{Les deux permutations du contour.} Le recollement identifie les deux côtés
d'une même arête de $\Gamma$ : c'est l'involution sans point fixe
\[
\sigma_{C} = \mathrm{rd} \circ \sigma_{\Gamma} \circ \mathrm{rd}^{-1}.
\]
La permutation « arête suivante » sur le contour orienté fait le tour d'une
face : au bout de l'arc $\bar{a}$, on repart par l'arc qui suit $\sigma_{\Gamma}
\bar{a}$ autour de son origine. D'où, en identifiant $A(C)$ à $\bar{A}$ par
$\mathrm{rd}$,
\[
\sigma_{C} = \sigma_{\Gamma}, \qquad
\rho_{C} = \rho_{\Gamma}\,\sigma_{\Gamma}, \qquad
\rho_{\Gamma} = \rho_{C}\,\sigma_{C},
\]
les cycles de $\rho_{C}$ étant les polygones du contour, c'est-à-dire les faces,
et les orbites de $\rho_{C}\sigma_{C}$ les sommets.\footnote{Selon qu'on longe
les arcs par la droite ou par la gauche et qu'on tourne dans un sens ou dans
l'autre, on trouve $\rho_{\Gamma}\sigma_{\Gamma}$ ou une variante conjuguée ; on
fixe la convention qui donne les deux dernières lignes du tableau de la page 9.
Les feuillets donnent le dictionnaire sous trois formes qui ne s'accordent pas.
La page 7 encadre « $\rho_{C}(\mathrm{rd}\,\bar{a}) =
\mathrm{rd}(\rho_{\Gamma}(\bar{a}))$ » et « $\rho_{C}(\mathrm{rg}\,\bar{a}) =
\mathrm{rg}(\rho_{\Gamma}(\bar{a}))$ », où $\sigma_{\Gamma}$ ne figure pas,
l'exposant de la seconde étant noyé d'encre. La page 9 pose en tête
« $\sigma_{\Gamma} = \sigma_{C}$, $\rho_{\Gamma} = \rho_{C}$ ». Enfin son tableau
associe $\sigma_{\Gamma}$ à $\rho_{C}$, $\rho_{\Gamma}\sigma_{\Gamma}$ à
$\rho_{C}$, et $\rho_{\Gamma}$ à $\rho_{C}\sigma_{C}$. Les deux dernières
correspondances sont celles qu'on écrit ; la première égalité de la tête aussi.
La page 7 construit encore des involutions $\sigma^{C}_{0}$, $\sigma^{C}_{1}$ sur
les arêtes orientées du contour, $A(C) \times \{\pm 1\} \simeq \bar{A} \times
\{\pm 1\}$, qui sont les opérations « changer de sommet » et « changer d'arête »
d'une carte de dimension 1 ; on ne les reprend pas.}

\textbf{Équivalence.} La catégorie des graphes combinatoires circularisés sans
sommet isolé est équivalente à celle des contours orientés $(A(C), \rho_{C})$
munis d'une involution sans point fixe $\sigma_{C}$ de $A(C)$. Dans les deux sens,
l'ensemble sous-jacent est le même, et l'on passe d'un couple de permutations à
l'autre par les formules ci-dessus.\footnote{C'est la description moderne d'une
\emph{carte combinatoire} par un ensemble de demi-arêtes muni d'une involution
et d'une rotation, les faces étant les cycles de leur produit. Le passage d'une
présentation à l'autre n'est qu'un changement de générateurs ; les noms ne sont
pas sur la page. Deux notes en biais dans la marge de la page 7 ne se lisent
pas.}

\textbf{Critère} (encadré, page 9). $\Gamma$ est un arbre, non ponctuel, si et
seulement si $C$ est un polygone connexe et que l'involution $\sigma_{C}$ ne
croise pas : pour deux orbites disjointes $\{a, \sigma_{C} a\}$ et
$\{b, \sigma_{C} b\}$, les points $b$ et $\sigma_{C}b$ sont du même côté de la
corde qui joint $a$ à $\sigma_{C}a$ dans l'ordre circulaire du polygone.

En effet, $\Gamma$ est connexe puisque $\rho_{C}$ n'a qu'un cycle, et la formule
d'Euler $|S| - |A| + 1 = 2 - 2g$ montre que $\Gamma$ est un arbre si et seulement
si la surface obtenue en recollant le polygone est une sphère, $g = 0$. Or le
recollement d'un polygone par un appariement de ses côtés, en renversant les
orientations, donne une sphère exactement quand l'appariement est sans
croisement.\footnote{La démonstration est la nôtre ; la page encadre l'énoncé,
dont les deux derniers mots, après « la composition est », ne se lisent pas. On
donne à « totalement non … » le sens que l'énoncé demande, celui d'un
appariement sans croisement.}

\pagerange{10}{12}
\section*{Disques découpés par cordes (pages 10 et 12)}

\emph{Retour au cas général.} Topologiquement, $|\Gamma|$ se déduit de $|C|$
par passage au quotient : on recolle chaque arête $a$ du contour avec l'arête
$\sigma_{C}(a)$, en renversant les orientations induites.\footnote{Il écrit
l'involution en exposant, $\sigma_{C}^{a}$.}

\emph{Le cas des arbres.} Quand $\Gamma$ est un arbre, $C$ est un seul polygone,
et l'on peut le voir comme le bord d'un disque orienté. En joignant par une
corde, à l'intérieur du disque, les milieux des arêtes qui se correspondent par
$\sigma_{C}$, on obtient des cordes qui ne se rencontrent pas, par le critère de la
page 9. Le disque est ainsi découpé en régions, et $\Gamma$ exprime la manière
dont les régions se recollent le long de leurs frontières communes : un sommet
par région, une arête par corde. C'est l'arbre dual du découpage.

\textbf{Énoncé final.} Il en résulte deux équivalences, en tête de la page 10
et à la page 12\footnote{« (non pointés) » et « (pas nécess.\ plus orientés) » sont ajoutés en
interligne. Pointés en un arc, les arbres plans à $n$ arêtes sont comptés par le
nombre de Catalan $\frac{1}{n+1}\binom{2n}{n}$, qui compte aussi les
appariements sans croisement des $2n$ côtés d'un polygone ; la remarque est la
nôtre. Le feuillet s'arrête sur l'énoncé.} :
\begin{enumerate}
\item[(i)] les arbres plans non pointés sont équivalents aux polygones
combinatoires munis d'une involution sans point fixe et sans croisement sur
l'ensemble de leurs arêtes ;
\item[(ii)] la catégorie des arbres plans est équivalente à la catégorie
isotypique des disques, non nécessairement orientés, munis d'un découpage par des
cordes qui ne se croisent pas ; les sommets de l'arbre correspondent aux régions,
c'est-à-dire aux paquets d'arêtes du polygone qui bordent une même région.
\end{enumerate}

\end{document}
